% 10-federation.tex: независимые органы.
\section{Федерация: независимые органы, признающие друг друга}
\label{sec:federation}

\subsection{Почему центра нет}

Возможны были две формы. Центральный экземпляр держит личность каждого человека и токены каждого
ведомства под одним корнем доверия, а ведомства суть его арендаторы. Федеративная форма даёт
каждому выдающему органу собственный экземпляр и собственный подписывающий ключ, записывает
межведомственное доверие явно и позволяет доверяющей стороне проверять против опубликованных
ключей без всякой центральной службы на пути. Проектная запись выбирает вторую и говорит, что выбор
этот не свободен: конституция называет граф доверия федерации несущим, гласит, что ни одно
ведомство не держит монополии, и отвергает сведение данных в масштабе населения структурно.
Центральный экземпляр есть монополия по построению и единое хранилище, которое можно захватить,
истребовать по повестке или перелопатить. Федеративная форма намеренно дороже; её цена есть
межведомственный протокол, и цена эта принята. \sImpl

\subsection{Доверие есть строка, а не цепочка}

Аттестация есть одна строка \entity{АттестацияДоверия}: \emph{проверяющее ведомство П
принимает выдающее ведомство Э в контексте К до такой-то даты}. Проверка между ведомствами удаётся
тогда и только тогда, когда ровно одна действующая строка отвечает верификатору, эмитенту и
контексту. Поиск есть один нерекурсивный запрос. Он никогда не вычисляет замыкания, так что доверие
П к Э никогда не влечёт доверия П к тому, чему доверяет Э; циклы возникнуть не могут, потому что
граф никто не обходит; а оператор, желающий доверия через несколько переходов, объявляет каждый
переход, так что след аудита держит каждое решение, а не посылки умозаключения. Самоаттестация
отвергается ограничением целостности. Отзыв задаёт дату и причину разом, единожды, и смотрит
вперёд: события, записанные, пока ребро было действующим, не обессмысливаются. \sImpl

Один предел назван, а не замазан. Строка записывает, какой оператор Полярис создал аттестацию; это
не криптографическая подпись аттестующего ведомства. Подписание аттестаций на уровне ведомства есть
миграция, под которую схема уже сформована, и задача управления ключами, которую книга готовности
возлагает на развёртывающую организацию. \sExt

\figfile{federation}{Два органа и одна направленная аттестация. Б аттестует ключ А в одном
контексте; доверяющая сторона, доверяющая Б, принимает удостоверение А в этом контексте, автономно,
из описи Б и ленты отзывов А. А доверяет также В, и Б это не даёт ничего: доверие никогда не есть
замыкание. В НИ это выполняется как два независимых экземпляра поверх ХТТП при смешанной паре
наборов параметров.}{fig:federation}

\subsection{Три подписанных объекта, которые публикует орган}

Доверяющей стороне нужны от чужого органа три вещи, и все три суть подписанные представления над
таблицами, работающими только на добавление, публикуемые на короткоживущих конечных точках и
потребляемые автономно отделяемым верификатором. Опись федерации несёт собственные якоря органа,
всякий ключ, каким он когда-либо пользовался, с его настоящим состоянием, и сделанные им
аттестации, каждую с аттестованным ключом. Контрольная точка эпохи есть его обязательство о
последней точке его цепи эпох. Лента отзывов есть множество отозванных листьев. Опись обязана быть
внутренне согласованной, подписанной одним из действующих якорей, которые она объявляет, так что
подписать её чужим ключом нельзя; подлинной, при двух свидетелях; свежей, в пределах окна, которое
потребитель ограничивает; и доверенной, от органа, чей якорь доверяющая сторона выбрала.
Межведомственное решение читается тогда так: принять тогда и только тогда, когда подпись
удостоверения подлинна, доверенная опись аттестует этот подписывающий ключ в предъявленном
контексте, а лента эмитента, падающая закрытой, удостоверения не перечисляет. \sImpl

Держатель способен также доказать принадлежность эпохе чужого органа с нулевым разглашением.
Доказательство уже несёт корень эпохи как открытый вход; корень уже едет по федерации подписанным
в контрольной точке. Решение складывает эти два: довериться чужой контрольной точке через граф,
чтобы получить доверенный корень, потребовать, чтобы доказательство было связано с этим корнем,
эпохой и контекстом, а затем проверить доказательство, вызвав двоичный файл доказывателя как
местный подпроцесс, без всякой сети. Если двоичного файла нет, верификатор воздерживается, что
ложным принятием не бывает никогда. Вердикт называет орган и решение, но никогда не токен. \sImpl

\subsection{Один экземпляр, много органов и та часть, которой политики не исправят}

Экземпляр Полярис способен держать много ведомств; схема это допускает, а поставляемый посев
показывает шесть. Ревью на \code{в9.364} спросило, способны ли операторы одного органа видеть
данные другого в таком экземпляре, и обнаружило, что способны: операторская учётная запись вообще
не несла привязки к органу, а шестнадцать из семидесяти четырёх операторских маршрутов читали
данные удостоверений или держателей без всякого фильтра по выдающему ведомству, и проще всех был
список токенов, отбиравший каждый токен, соединённый с каждым держателем, без всякого понятия о
том, чей он. Операторы теперь привязаны к ведомству, и каждый такой маршрут по нему фильтрует. Что
ревью записало ещё, так это ту часть, которой никакой фильтр не исправит: в общем экземпляре
владелец базы данных, резервная копия и схема суть общие, так что изоляция внутри одного экземпляра
есть любезность, оказываемая приложением, а не граница, которую держит схема. Топология,
выбранная проектной записью, состоит в том, что каждый орган запускает собственный экземпляр, а
общий экземпляр есть решение развёртывания, цена которого изложена. \sImpl{} для привязки и
фильтров; \sExt{} для границы между органами, которая есть развёртывание.

\subsection{Два экземпляра и что они доказывают}

Отдельная задача НИ поднимает два независимых экземпляра как два органа, у каждого своя база
данных и свой настоящий корень МЛ-ДСА, говорящих только поверх ХТТП, и гоняет всю матрицу:
межведомственную проверку по забранной описи, отзыв аттестации, переворачивающий решение,
контрольные точки эпох и ленты отзывов, забранные и проверенные по проводу, компрометацию ключа,
записанную на одной стороне и обращающую принятие на другой, шлюз обмена, службу отметок времени,
прочитанную из забранного реестра, вход через брокер, подписание документа через кошелёк и
смешанную пару наборов параметров (МЛ-ДСА-87 с одной стороны, МЛ-ДСА-65 с другой). Она краснеет
на любом неверном решении. \sImpl

Что эта задача доказывает, так это что протокол работает между двумя экземплярами этой же кодовой
базы, запущенными одним автором на одной машине. Она не доказывает институциональной
независимости и не доказывает, что независимая реализация, написанная кем-то другим по
спецификации, взаимодействует по родному протоколу. Первое требует развёртывания учреждениями,
которые не есть автор. Второе требует сторонней команды и не случилось; единственная сторонняя
реализация, заговорившая с Полярис, сделала это по ОИД4ВП (раздел~\ref{sec:verifiers}), а не по
этим форматам. Репозиторий проводит эту черту сам: независимость протокола, независимо
эксплуатируемые учреждения и независимые сторонние реализации суть три разных утверждения, и
показано лишь первое. \sExt

\begin{layercard}{Карточка слоя: федерация}
\qa{Что это}{Корни у каждого органа, направленные поконтекстные аттестации, подписанные опись, контрольная точка эпохи и лента отзывов, сводный набор состояний, межведомственная принадлежность с нулевым разглашением, изоляция операторов по органам и учение на двух экземплярах.}
\qa{Зачем оно существует}{Ни одно ведомство не вправе держать монополию, и не должно существовать единого хранилища всех людей; и всё же доверяющая сторона обязана уметь принять удостоверение, выданное в другом месте.}
\qa{Чему доверяет}{Якорям, которые доверяющая сторона выбрала сама; таблицам только на добавление, над которыми подписанные объекты суть представления.}
\qa{Чему отказывает}{Транзитивному доверию; центральной службе на пути проверки; описи, не подписанной её собственным объявленным якорем; слову сводчика о состоянии участника; обратному вызову, который эмитент мог бы записать; оператору, читающему держателей другого органа.}
\qa{Что видит}{Опубликованные данные доверия: ключи, аттестации, корни эпох, отозванные листья. Ни одного человека.}
\qa{Чего не видит}{Действующее население любого органа; кто кого проверял.}
\qa{Если откажет}{Поддельная опись падает на внутренней согласованности; раздвоенная цепь эпох или откаченная лента суть неотрицаемое свидетельство; скомпрометированный ключ эмитента выводится из обращения с указанного момента списком доверия, которым эмитент не распоряжается; владелец общего экземпляра лежит за пределами того, что способен ограничить любой фильтр.}
\qa{Кем проверяется}{\code{проверка\_топология\_федерации}, \code{проверка\_межведомственный\_протокол}, \code{проверка\_распространение\_отзывов\_по\_эпохам}, \code{проверка\_межведомственное\_нр}, \code{проверка\_федерация\_два\_экземпляра}, \code{проверка\_доверие\_обмена\_направленное}, \code{проверка\_стражи\_маршрутов\_выводятся\_а\_не\_перечисляются}.}
\qa{Сейчас или потом}{Сейчас, как протокол между экземплярами. Независимые учреждения и реализации потом, и поставлять их не нам.}
\qa{Внутри и снаружи}{Внутри: использование. Снаружи: прозрачность, потому что подписанную историю всё ещё можно рассказать разным людям по-разному.}
\end{layercard}

\nextlayer{Подписанные объекты федерации доказывают, что орган сказал. Нечестный орган всё ещё
способен говорить разное разным наблюдателям, публикуя одну цепь эпох аудитору и другую доверяющей
стороне. Поэтому Полярис нужна прозрачность: способ для наблюдателей доказать, что история только
росла, и сверить записи между собой.}
